\honest{37 Ways That Words Can Be Wrong}{Eliezer Yudkowsky}{2008}

Some reader, I expect, will say this post should be called ``37 Ways That You Can Use Words Unwisely,'' or something longer about suboptimal categories. My answer is that saying ``There's no way my choice of X can be `wrong'\,'' is ``nearly always an error in practice, whatever the theory.'' You can always be wrong; there is never a get-out-of-jail-free card. Besides, I can define ``wrong'' any way I like.

Here is the list. Each item names an error, gives an example, and points to the post where I discussed it. The first 36 items cover every earlier post of the sequence, in order. \nb{Criticisms of those posts are in the notes on each of them. The notes here concern only the summary.}

First, words that fail. A word can connect to nothing: is Socrates a ``framster''? An argument can make reality depend on a definition, or claim an empirical fact is true ``by definition,'' though ``Logical truths are true in all possible worlds, and so never tell you which possible world you live in.'' You define ``human'' by mortality and still call Bob human without watching him die. Calling blue eggs ``bleggs'' hides an inference from eleven eggs and eight cubes. You define ``red'' with words and never point to a stop sign.

Then the Mars item: ``We aren't consciously aware of our identification of a red light in the sky as `Mars'.'' \nb{``Extensions and Intensions'' had said something different: that we are not aware that this identification ``is a separate matter from our verbal definition.''}

Next come categories. Definitions capture only part of a cluster, as Diogenes's plucked chicken showed. Membership is graded: people think a disease spreads more easily from robins to ducks than the reverse. Nitpicking a definition that works is an error too. Whether Barney the Baby Boy is a ``man'' depends on whether you want to know if hemlock is good to feed him or if he will make a good husband. Human brains notice blegg against rube far more easily than that red objects never glow though red furred ones are otherwise bleggs, and ``other forms of statistical inference are possible even though your brain doesn't use them.'' \nb{``Neural Categories'' had offered the claim about the brain as ``a fair guess.''} Categories are inferences in a brain, not gifts from Plato. Once every property of an object is known, nothing is left to ask by arguing whether it is a blegg.

Then arguments over words. Disputes slide into definitions. Words feel as if their meaning were part of them. People argue over meanings after both sides understand each other. Dictionaries record usage; do not pull one out in an empirical or moral argument, or ``in the middle of any argument ever.'' \nb{The source post limited this to arguments where an empirical or moral proposition is at stake, and the very next item counts defying common usage as an error. In a dispute about usage itself, a dictionary is a record of usage.} Do not defy common usage without a reason (``Fast stand up plutonium, with bagels without handle''). Renaming creates the illusion of inference. Banning a word, as in Taboo, dissolves some arguments. A neat word hides details (what is ``truth'' if you cannot say ``accurate,'' ``correct,'' ``map'' or ``real''?), and one word can hide two things, as when a detective fails to see there are two Carols. Labels breed stereotypes, as with blood types in Japan. People sneak in connotations, as with the ``wiggin,'' and say ``by definition'' when a default inference has been called into doubt.

Then boundaries. A definition can group things that do not belong together, like dolphins among fish. Short words should go to frequent things; otherwise your mind may think short sentences sound ``simpler,'' as ``God did a miracle'' does. \nb{The cited post, ``Entropy, and Short Codes,'' does not make this point. The earlier post ``Occam's Razor'' does: ``the length of an English sentence is not a good way to measure `complexity.'\,''} A word for a region where nothing clusters, like ``wiggin'' for green eyes and black hair, allows no inferences. Drawing an odd boundary without a reason, such as a word for all humans except black people, is like a detective naming a suspect with no evidence. One item warns against inferring properties by category when they are not independent given the class, as Naive Bayes assumes; that one I do not try to summarize: ``Just read the blog post.''

Last, words and minds. Words are not ``tiny little LISP symbols'' but handles for mental paintbrushes. A word can mean different things in different places, like ``left'' in ``Martin told Bob the building was on his left.'' And the thirty-seventh: thinking that definitions cannot be wrong, which ``teaches you to indignantly defend your past actions.'' For its source I cite this post.

Your brain races ahead without your supervision. Saying ``Words are arbitrary; I can define a word any way I like'' is like driving on thin ice, accelerator floored, and turning the wheel at random because no angle looks special. \nb{The sequence had separated these two claims. ``The Argument from Common Usage'' grants that ``we have no obvious mutual interest in using any particular word for a concept.'' Where a concept's boundary falls is what the sequence argued can be mistaken. The steering wheel fits only that second claim.} Pay attention to ``the three or six dozen optimality criteria'' that govern how you use words.
